%% RESUMO NA LINGUA VERNACULA (elemento obrigatorio -- NBR 6028)
%% Paragrafo unico, sem recuo, 150 a 500 palavras; palavras-chave em
%% config/dados.tex. Versao provisoria (dissertacao dummy).

\begin{resumo}
Acervos educacionais audiovisuais mantidos por redes públicas de ensino
constituem ativos digitais do Estado, cujo valor público depende menos de sua
existência do que da capacidade institucional de convertê-los em uso pedagógico
efetivo. Esta pesquisa desenvolve e avalia um artefato de recomendação
pedagógica baseado em inteligência artificial que classifica as 760 videoaulas
ativas do Preparatório ENEM do Canal Educação, programa da Secretaria de Estado
da Educação do Piauí, segundo as 120 habilidades da matriz de referência do
Exame Nacional do Ensino Médio, com o propósito de apoiar o planejamento docente
sem substituir a decisão pedagógica. Adota-se a Design Science Research como
método. O artefato converte o conteúdo audiovisual em texto por transcrição
automática, representa aulas e descritores por embeddings do modelo multilíngue
BGE-M3, empregado exclusivamente em inferência, e ordena, para cada aula, as
habilidades mais próximas pela similaridade de cosseno, sem depender de
histórico de uso. A avaliação técnica confronta o método do artefato, um método
lexical sem inteligência artificial (TF-IDF) e um modelo generativo com a
atribuição oficial de habilidades a itens do exame publicada pelo INEP. Sobre os
178 itens da edição de 2023, o artefato supera o método lexical na primeira
posição do ordenamento e em cinco posições, mas alcança menos da metade da
concordância do modelo generativo. Os resultados indicam a viabilidade técnica
da recomendação baseada em conteúdo em contexto de início frio e evidenciam
limites que condicionam a interpretação das saídas. Do ponto de vista da gestão
pública, a classificação do acervo converte-se em insumo de governança do
próprio ativo, ao expor as lacunas de cobertura curricular que orientam a
priorização da produção de novos conteúdos.
\end{resumo}
